\chapter{State of the Art}\label{ch:sota}

The state of the art in cache-aware trie index structures falls into six thematic
search clusters (A--F) and a parallel allocator corpus (A01--A23). In total, 33
search-algorithm papers and 21 allocation papers were studied in depth; the results are
available as persisted XML configurations under \texttt{algorithm\_profiles/} in the
\texttt{comdare-cache-engine}. The complete building-block and allocator matrices are
given in Appendix~\ref{app:blocks}.

\section{Trie Family (Cluster A)}\label{sec:cluster-a}

\textbf{P01 ART}~\cite{leis2013art} introduces the Adaptive Radix Tree with adaptive node
sizes Node4/16/48/256. \textbf{P02 HOT}~\cite{binna2018hot} optimizes trie height via
multi-byte prefixes and bit-vector indexes. \textbf{P03 Masstree}~\cite{mao2012masstree}
combines a trie with B-trees and a per-internal-node permutation index. \textbf{P04
CoCo-Trie}~\cite{boffa2024coco} uses succinct-encoded compact pages. \textbf{P05
START}~\cite{fent2020start} extends ART with self-tuning of the node representation. P09
LOUDS (Jacobson 1989) is the succinct-trie foundation, and \textbf{P10
SuRF}~\cite{zhang2018surf} uses LOUDS for range filtering.

\section{Hybrid and B\texorpdfstring{\textsuperscript{+}}{+} Family (Cluster B)}\label{sec:cluster-b}

\textbf{P06 B\textsuperscript{2}-tree}~\cite{schmeisser2022b2tree} employs a two-level
fanout for combined L1+L2 cache awareness. \textbf{P07 Wormhole}~\cite{wu2019wormhole}
combines a hash table with a linear prefix view. P11 CSS-tree (Rao/Ross 1999) is the first
cache-sensitive search tree; P12 CSB\textsuperscript{+}-tree (Rao/Ross 2000) develops it
into a B\textsuperscript{+}-tree; P13 Hankins/Patel (2003) studies the effect of node size
on cache performance.

\section{Layout Theory (Cluster C)}\label{sec:cluster-c}

P14 Samuel/Pedersen/Bonnet (2005) makes the CSB\textsuperscript{+}-tree
processor-conscious. P15 Graefe/Larson (2001) studies B-tree indexes versus CPU caches.
P16 Bender/Demaine/Farach-Colton (2002) introduces cache-oblivious tree layout (van Emde
Boas); P17 Bender et al.\ (2005) extends it to cache-oblivious B-trees. P18
Saikkonen/Soisalon-Soininen (2008) and P19 (2016) develop multi-level and
layout-invariant B-trees.

\section{Prefetching (Clusters D and E)}\label{sec:cluster-de}

P20 ``B-Trees Are Back'' (Mueller/Benson/Leis 2025) demonstrates modern B-tree
engineering. P21 Chen/Gibbons/Mowry (2001) introduces prefetching B\textsuperscript{+}-trees;
P22 Chen et al.\ (2002) develops fractal prefetching. P23 Khan (2010) makes cache
prefetching dynamically adaptive; P24 Naderan-Tahan/Sarbazi-Azad (2016) adds an adaptive
filter on the cache prefetcher. P25 Mahling/Weisgut/Rabl (2025) combines range scans with
hot-path caching. P26 Zhang (FGCS 2024) and P27 Zhang (ASPLOS 2025, hierarchical
prefetching) are the most recent schemes, and P28 Kuehn (DaMoN 2023) motivates
data-driven cache optimization.

\section{Synchronization (Cluster F)}\label{sec:cluster-f}

P08 ``ART of Practical Synchronization'' (Leis 2016) introduces optimistic lock coupling
(OLC) for ART. P29 RCU (McKenney, OLS 2001) and P30 Hazard Pointers (Michael 2004) are the
canonical reclamation schemes. P31 Ungethuem/Habich (2017) and P32 Schmidt/Habich (2025)
are TU-Dresden-specific hardware-optimization surveys; P33 (VAMPIR poster,
Berthold/Habich 2023) provides the OTF2 profiler integration.

\section{Allocator Research (A01--A23)}\label{sec:allocators}

The 21 allocation papers are catalogued systematically in the allocator matrix
(Appendix~\ref{app:blocks}). Tier-1 representatives are Hoard (A01, ASPLOS 2000), Slab
(A02, USENIX 1994), Michael lock-free (A03, PLDI 2004), mimalloc (A04, MSR-TR-2019-18),
jemalloc (A05), TCMalloc (A06), snmalloc (A07, ISMM 2019), and NUMAlloc (A09, ISMM 2023).
Tier-2/3 add CAMA (A12, RTAS 2011), StarMalloc (A13, OOPSLA 2024), Crystalline reclamation
(A17, PLDI 2024), and PIM-malloc (A16, arXiv 2025).

\section{Profile Schema with the \texorpdfstring{\texttt{<expected\_workload>}}{expected\_workload} Tag}\label{sec:profile-schema}

For each SOTA algorithm and each allocator, a profile XML in
\texttt{cache\_engine/algorithm\_profiles/} records metadata, the axis assignment (page,
node, traversal, value\_handle, concurrency, allocator, prefetch, telemetry, isa, layout,
reclamation), a key--value signature, and an optional \texttt{<expected\_workload>} tag
that overrides the \texttt{ExperimentDriver} heuristic (Chapter~\ref{ch:methodology}). All
30 SOTA profiles and all 10 allocator profiles are tagged.

\begin{table}[!htbp]
\caption{SOTA algorithm profiles with workload affinity}\label{tab:sota-profiles}
\centering\small
\begin{tabular}{lll}
\toprule
\textbf{Profile} & \textbf{Paper ref.} & \textbf{expected\_workload} \\
\midrule
art & P01 (Leis 2013) & YCSB\_C \\
hot & P02 (Binna 2018) & YCSB\_C \\
masstree & P03 (Mao 2012) & YCSB\_A \\
coco\_trie & P04 (Boffa 2024) & YCSB\_C \\
start & P05 (Fent 2020) & YCSB\_C \\
b2tree & P06 (Schmeisser 2022) & YCSB\_E \\
wormhole & P07 (Wu 2019) & YCSB\_A \\
surf & P10 (Zhang 2018) & YCSB\_A \\
css\_tree & P11 (Rao/Ross 1999) & YCSB\_C \\
csb\_tree & P12 (Rao/Ross 2000) & YCSB\_E \\
hankins & P13 (Hankins 2003) & YCSB\_C \\
samuel & P14 (Samuel 2005) & YCSB\_C \\
graefe & P15 (Graefe 2001) & YCSB\_C \\
bender\_treelayout & P16 (Bender 2002) & YCSB\_C \\
bender\_cacheoblivious & P17 (Bender 2005) & YCSB\_C \\
saikkonen\_multilevel & P18 (Saikkonen 2008) & YCSB\_E \\
saikkonen\_layoutinvariant & P19 (Saikkonen 2016) & YCSB\_E \\
btreesareback & P20 (Mueller 2025) & YCSB\_E \\
chen\_prefetch & P21 (Chen 2001) & YCSB\_A \\
chen\_fractal & P22 (Chen 2002) & YCSB\_A \\
khan\_adaptive & P23 (Khan 2010) & YCSB\_A \\
naderan\_tahan & P24 (Naderan-Tahan 2016) & YCSB\_A \\
mahling & P25 (Mahling 2025) & YCSB\_E \\
zhang\_fgcs & P26 (Zhang 2024) & YCSB\_A \\
zhang\_asplos & P27 (Zhang 2025) & YCSB\_A \\
kuehn & P28 (Kuehn 2023) & YCSB\_C \\
rcu & P29 (McKenney 2001) & YCSB\_B \\
hazard\_pointers & P30 (Michael 2004) & YCSB\_B \\
ungethuem & P31 (Ungethuem 2017) & YCSB\_C \\
to\_stride & P32 (Schmidt 2025) & YCSB\_C \\
\bottomrule
\end{tabular}
\end{table}

The workload distribution is 13$\times$ YCSB\_C (read-heavy), 9$\times$ YCSB\_A (Zipfian
read+update), 6$\times$ YCSB\_E (range scans), and 2$\times$ YCSB\_B (95/5 read/update).
PRT-ART itself is tagged YCSB\_F (read--modify--write, for density-tracker updates).
Table~\ref{tab:allocator-profiles} lists the ten implemented allocator profiles.

\begin{table}[!htbp]
\caption{Allocator profiles with license and workload affinity}\label{tab:allocator-profiles}
\centering\small
\begin{tabular}{llll}
\toprule
\textbf{Profile} & \textbf{Family} & \textbf{License} & \textbf{expected\_workload} \\
\midrule
hoard & A01 & Apache-2.0 & YCSB\_A \\
michael\_lockfree & A03 & BSD-3 & YCSB\_B \\
mimalloc & A04 & MIT & YCSB\_A \\
jemalloc & A05 & BSD-2 & YCSB\_B \\
tcmalloc & A06 & BSD-3 & YCSB\_C \\
snmalloc & A07 & MIT & YCSB\_A \\
scalloc & A08 & BSD-3 & YCSB\_A \\
rpmalloc & A10 & Public Domain & YCSB\_C \\
lrmalloc & A11 & BSD-3 & YCSB\_B \\
dlmalloc & A20 & CC0 & YCSB\_C \\
\bottomrule
\end{tabular}
\end{table}

\section{Hardware and Scheduling Strategy per Paper}\label{sec:hw-sched}

After the N-phase extension of the building-block matrix (Section~\ref{sec:axes}), two
additional algorithm permutation axes are anchored:

\begin{itemize}
  \item \textbf{Axis~12 Hardware-Strategy}: which hardware features the algorithm actively
        uses --- sub-axes 12.1 SIMD family (AVX2/AVX-512/NEON/SVE2/scalar), 12.2
        cache-level targeting (L1/L2/L3/HBM-aware), 12.3 NUMA strategy, 12.4 prefetch
        hardware (PREFETCH/PREFETCHNTA/PREFETCHW), 12.5 atomic-instruction family
        (CAS/LL-SC/RMW-extended).
  \item \textbf{Axis~13 Scheduling-Strategy}: how work is distributed --- sub-axes 13.1
        worker-pool layout (thread-per-core / work-stealing / CPU-pinning), 13.2
        SIMD-worker-count limit (hardware-limited to typically 2 of $N$ cores), 13.3
        heterogeneous-core dispatch (Intel hybrid P-/E-cores), 13.4 co-routine strategy,
        13.5 batch granularity.
\end{itemize}

Table~\ref{tab:hw-sched} shows a selection of these two axes per SOTA paper; the full
catalogue is in Appendix~\ref{app:blocks}.

\begin{table}[h!]
\centering\small
\caption{Hardware and scheduling strategy per SOTA paper (selection)}\label{tab:hw-sched}
\begin{tabular}{lll}
\toprule
\textbf{Paper} & \textbf{Hardware strategy (axis 12)} & \textbf{Scheduling strategy (axis 13)} \\
\midrule
P01 ART & AVX2, L1-aware, none, CAS & thread-per-core \\
P02 HOT & AVX2 (bit-vector), L1-aware, CAS & thread-per-core \\
P03 Masstree & scalar, L1-aware, CAS & thread-per-core + work-stealing \\
P05 START & AVX2 (multi-byte span), L1-aware, CAS & thread-per-core \\
P20 LeanStore & AVX2, HBM-aware, NUMA-aware, CAS & work-stealing + hybrid \\
P27 hp-soft & PREFETCH (L1/L2/L3 bundle), all-cache, CAS & thread-per-core + co-routine \\
P28 Kuehn & PREFETCH (scalar counter), L1-aware, CAS & thread-per-core + sampling \\
P31 Ungethuem & AVX-512 (HBM), HBM-aware, NUMA-aware, CAS & hybrid (P-/E-cores) \\
P32 Schmidt & AVX-512 (SIMD stride), HBM-aware, NUMA-aware, CAS & hybrid \\
\midrule
\textbf{PRT-ART} & \textbf{AVX2, L1-aware (TLB-aligned), local NUMA, PREFETCH, CAS (OLC)} & \textbf{thread-per-core + 2-SIMD-worker limit + hybrid + co-routine + micro-batch} \\
\bottomrule
\end{tabular}
\end{table}

\section{Research Gap}\label{sec:gap}

Each of the surveyed works optimizes one or a few axes of the cache hierarchy in
isolation, and the results are reported in incomparable settings. What is missing --- and
what this thesis contributes --- is a systematic, orthogonal \emph{axis decomposition} of
these designs together with a permuted measurement architecture that reconstructs them as
explicit configurations and measures them on one common basis.
